\section{INTRODUCCIÓN}

 Detectar phishing con reglas fijas tiene un problema de fondo y es que las reglas se pueden aprender y evadir. Por eso gran parte de la investigación de la última década se centró en el aprendizaje automático, que en lugar de comparar contra una lista aprende patrones a partir de miles de ejemplos etiquetados. Los primeros enfoques de este tipo se apoyaron casi por completo en características léxicas extraídas directamente de la URL como la longitud, número de puntosy presencia de palabras sospechosas, combinadas con clasificadores clásicos como Random Forest o SVM, y llegaron a superar el 97\% de exactitud usando conjuntos de datos de más de 70.000 URLs \cite{sahingoz2019machine}. Otros trabajos sumaron a esas características léxicas información obtenida del código fuente y de servicios de terceros como la edad del dominio, certificados o reputación, lo que llevó la exactitud hasta 99,31\% con Random Forest \cite{rao2019detection}. El problema de tener demasiadas variables, muchas de ellas redundantes, motivó el desarrollo de esquemas de selección de características por ensamble, donde uno de ellos logró distinguir el 94,6\% de los sitios usando solo el 20,8\% de las características originales \cite{chiew2019hybrid}, mientras que otro combinó esa selección con una red neuronal para obtener un modelo más compacto sin perder desempeño \cite{zhu2019ofsnn}. En paralelo surgieron modelos que evitan por completo la ingeniería manual de características, procesando la URL carácter por carácter mediante redes convolucionales \cite{aljofey2020effective}. \newline

Sobre esa base, la investigación más reciente se movió hacia arquitecturas de apilamiento con modelos de apilamiento en múltiples capas, que combinan clasificadores como MLP, KNN, Random Forest, Regresión Logística y XGBoost organizados en capas donde la salida de una alimenta a la siguiente, estas arquitecturas alcanzaron entre 96,79\% y 98,90\% de exactitud sobre distintos datasets de phishing \cite{kalabarige2022multilayer}, y una versión posterior de ese mismo enfoque, pero reemplazando los clasificadores base por algoritmos de boosting y sumando una etapa de selección de características por importancia promedio, mostró que reducir las variables a un 66-86\% del total no perjudica el desempeño, e incluso lo mejora en algunos casos \cite{kalabarige2023boosting}. Así mismo se ha explorado mantener fijo un solo algoritmo base y variar la estrategia de meta aprendizaje (AdaBoost, Bagging, Rotation Forest, LogitBoost), lo que también supera el 97\% de exactitud con tasas de falsos positivos por debajo de 0,03 y deja claro que tanto el clasificador base como la técnica de ensamble influyen en el resultado final \cite{alsariera2020ai}. Sin embargo, la exactitud reportada en estos trabajos no cuenta toda la historia, pues se ha demostrado que modelos como SVM caen a una tasa de detección de apenas 10\% frente a URLs generadas deliberadamente para evadirlos, mientras que los enfoques multimodales, que combinan dominio, HTML y contenido visual, resisten mucho mejor, con detección cercana al 99\% \cite{duy2024study}, y una revisión reciente de más de 130 artículos publicados entre 2020 y 2024 señala que buena parte de los modelos actuales siguen entrenados sobre datasets fijos y con extracción de características poco dinámica, lo que limita su capacidad de adaptarse a técnicas de phishing nuevas \cite{ahmad2025across}. En conjunto, estos trabajos confirman que los ensambles y el stacking superan a los clasificadores individuales, que la selección de características puede reducir el costo computacional sin sacrificar exactitud, y que ningún modelo está exento de ser evadido si no se piensa también en su robustez frente a ataques adversarios.